\chapter*{前言}
\addcontentsline{toc}{chapter}{前言}
\markboth{前言}{前言}

尽管\textbf{关联分析（association analysis）}在解析复杂疾病的遗传结构方面取得了显著进展，我们对于复杂疾病病因及其机制的认识仍然有限。关联分析中的显著发现往往缺乏一致性，也经常引发争议。目前，基因组分析无论在广度——一次分析的变量数量——还是在深度——遗传变异跨越基因组和分子层级、最终到达临床结局所经历的步骤数——方面，都存在不足；其分析范式仍以关联和相关分析为主。新一代基因组、表观基因组、传感与影像技术，正在不断产生层次更深入的多组学、生理、影像、环境与表型数据。对这些数据进行\textbf{因果推断（causal inference）}，是科学发现的基石，也是揭示疾病机制不可缺少的组成部分。为了阐明精神障碍的机制，现在有必要推动遗传分析范式的转变：从浅层关联分析转向深层因果推断，从单独开展遗传分析转向基因组、表观基因组、影像与表型数据的整合分析。

本书是《Big Data in Omics and Imaging: Association Analysis》的自然延伸，重点讨论基因组、表观基因组与影像数据的整合分析及因果推断。为使上述范式转变成为可能，本书将：
\begin{enumerate}[label=（\arabic*）]
\item 开发新的因果推断方法，并应用已有方法，开展复杂疾病的全基因组和全表观基因组因果研究。
\item 建立统一框架，对整合的基因组、表观基因组、影像与临床表型数据进行系统因果分析，推断多层级组学与影像因果网络，从而发现遗传变异经由这些网络通向疾病的路径。
\item 开发并应用基因表达与甲基化解卷积方法，发展新的细胞特异性多组学因果网络推断方法。
\item 介绍深度学习在基因组、表观基因组与影像数据分析中的应用，并发展将深度学习与因果推断相结合的方法。
\end{enumerate}

本书分为七章，各章内容如下。

第1章“基因型—表型网络分析”研究有向和无向的基因型—表型网络，这是因果推断的重要主题。有效的遗传分析包括两个主要方面：一是广度，即遗传变异影响的表型数量；二是深度，即遗传变异到达临床结局所经历的步骤数。因果推断理论与链图模型，为深入、精确地分析多层级混合型基因型—疾病因果网络，提供了新的分析平台。很少有遗传学与表观遗传学教材深入介绍因果推断理论，因此，第1章和第2章将为基因组与表观基因组分析中的因果推断，提供扎实的知识基础和有效工具。第1章内容包括：（1）用于基因型网络的无向图；（2）用交替方向乘子法估计高斯图模型；（3）坐标下降算法与图套索；（4）多个图模型；（5）网络的有向图与结构方程模型；（6）稀疏线性结构方程；（7）利用下一代测序数据构建基因型—表型网络的函数型结构方程模型；（8）效应分解与估计。

第2章“因果分析与网络生物学”介绍：（1）以贝叶斯网络作为因果推断的一般框架；（2）连续型因果网络的结构方程与评分度量；（3）用于学习混合贝叶斯网络的网络惩罚逻辑回归；（4）基于家系的因果推断统计方法；（5）非线性结构方程模型；（6）线性与非线性混合结构方程模型；（7）联合使用干预数据和观测数据进行因果推断；（8）用于因果结构学习的整数规划。

第3章“可穿戴计算与函数值性状的遗传分析”研究函数值性状的遗传学。疾病的早期发现与健康监测，是医疗保健和疾病管理的主要目标。ECG、EEG、SCG、EMG、MEG及血氧饱和度等生理性状，能够提供有关人体健康状态的重要信息，可用于疾病监测与诊断。可穿戴传感器能够以无创、连续的方式开展个体健康监测，不仅可测量静息状态下的健康参数，还能记录短暂事件所产生的信号；这些信号可能具有重要的预后或治疗意义。上述生理性状属于函数值性状。基因组数据与时空生理数据的联合分析，可以呈现疾病的整体遗传结构，同时也带来了很大的方法学与计算挑战。现有文献中，针对函数值性状开展遗传分析的统计方法仍较为缺乏，本章因此提出一些新的统计方法。内容包括：用于自动疾病诊断和实时健康监测的可穿戴计算；生理时间序列数据分析中的深度学习；响应变量与预测变量均为函数的函数型线性模型，用于结合下一代测序数据分析生理性状的关联；具有函数型响应的混合函数型线性模型，用于基于家系的生理性状遗传分析；响应变量与预测变量均为函数的函数型回归模型，用于基因—基因交互作用分析；以及用于生理性状关联研究的函数型典型相关分析。

第4章“RNA测序数据分析”介绍：（1）数据标准化与预处理；（2）针对RNA测序或miRNA测序数据开展差异表达分析的函数主成分分析检验；（3）用于等位基因特异性表达分析的多元函数主成分分析；（4）RNA测序数据的eQTL及eQTL上位性分析；（5）共表达网络；（6）线性与非线性调控网络；（7）基因表达插补；（8）基因型—表达调控网络；（9）动态贝叶斯网络与纵向表达数据分析；（10）单细胞RNA测序数据分析、基因表达解卷积与遗传筛查。

第5章“甲基化数据分析”讨论甲基化数据的分析。差异基因表达、eQTL分析和基因型—表达调控网络的统计方法，很容易扩展到甲基化数据分析。本章将首次引入表观遗传分析中的一个新概念：全表观基因组因果研究。除此之外，本章重点讨论全基因组甲基化与表达因果网络的推断。基因表达与甲基化数据都涉及两万多个基因，因此不可能直接构建一个包含四万多个节点的因果网络，需要设计多层级甲基化—表达网络。本章围绕这类网络估计中的三个关键问题展开：（1）用低秩模型表示某条通路或某个聚类中的基因表达或甲基化；（2）利用各通路或聚类中甲基化与基因表达的低秩表示，构建甲基化与表达网络；（3）在局部相连的通路或聚类中，利用原始甲基化与基因表达值，构建甲基化和基因表达因果网络。本章还研究哪些细胞中的甲基化会调控哪些细胞中的基因表达，并提出若干新的分析方法。

第6章“影像与基因组学”侧重于影像信号处理、自动影像诊断以及遗传—影像数据分析。为了开展疾病诊断、管理与机制研究，如何用统计方法和计算算法分析高维、具有空间相关性且结构复杂的影像数据，以及临床和遗传数据，正受到越来越多的关注。本章介绍：（1）用于医学影像语义分割的深度学习；（2）用于提取影像信号的三维函数主成分分析；（3）影像网络构建与连接性分析；（4）用于自动影像疾病诊断的因果机器学习；（5）利用下一代测序数据开展影像遗传分析的多元函数型线性模型；（6）用于影像遗传学或影像—RNA测序数据分析的二次正则化函数型典型相关分析；（7）影像遗传学及影像—RNA测序数据的因果分析；（8）用于整合fMRI与基因组数据开展因果分析的时间序列结构方程模型；（9）因果机器学习。

第7章“从关联分析到整合因果推断”发展新的全基因组因果研究统计方法，并研究基因组、表观基因组、影像与多表型数据的整合分析。本章给出因果分析的数学表述，讨论因果关系背后的原理，并介绍区分因果检验与关联检验的准则。在基因组和表观基因组数据分析中，通常考虑四类关联：离散变量与连续变量的关联、连续变量与连续变量的关联、离散变量与二元性状的关联，以及连续变量与二元性状（疾病状态）的关联。本章将这四类关联分析扩展为四类因果分析，还介绍若干强有力的因果推断工具，包括加性噪声模型、信息几何、迹方法、Haar测度与距离相关。

基因与表型之间存在多个环节。只有以足够的广度与深度，探索连接遗传变异和临床结局的庞大路径空间，才能揭示疾病机制。精准医学要求对基因型—表型开展深入、系统、全面而精确的分析——“探索得越深，了解得越多”。第7章提出，利用因果推断理论建立新的分析平台，对多层级混合型基因型—疾病因果网络开展深入而精确的分析。该平台把基因关联子网络、环境子网络、基因调控子网络、遗传—甲基化因果子网络、甲基化—基因表达网络、基因型—基因表达—影像子网络、中间表型子网络和多种疾病子网络，整合为一个相互连通的多层级基因型—疾病网络，以揭示疾病机制中的深层因果链。本章还讨论存在混杂因子时的因果推断。

总体而言，本书介绍因果推断、深度学习，以及基因组、表观基因组、影像和多表型数据分析中的前沿研究与实践成果，为这个充满挑战、迅速发展的领域后续研究奠定基础，并提供分析平台。我们希望，书中介绍的概念、统计方法、计算算法和分析平台，能够有助于培养新一代统计遗传学家、生物信息学家与计算生物学家。

感谢Sara A. Barton对本书的编辑工作。衷心感谢与我合作多年的同事和合作者Li Jin、Eric Boerwinkle等人。尤其感谢我过去和现在的学生及博士后，他们对研究的投入和科学贡献使本书得以完成：Jinying Zhao、Li Luo、Shenying Fang、Nan Lin、Rong Jiao、Zixin Hu、Panpan Wang、Kelin Xu、Dan Xie、Xiangzhong Fang、Jun Li、Shicheng Guo、Shengjun Hong、Pengfei Hu、Tao Xu、Wenjia Peng、Xuesen Wu、Yun Zhu、Dung-Yang Lee、Lerong Li、Getie A. Zewdie、Long Ma、Hua Dong、Futao Zhang和Hoicheong Siu。最后，感谢编辑David Grubbs，在本书写作过程中给予我的鼓励与耐心。

\medskip
MATLAB® 是The MathWorks, Inc.的注册商标。有关产品信息，请联系：
\begin{flushleft}
The MathWorks, Inc.\\
3 Apple Hill Drive\\
Natick, MA 01760-2098 USA\\
电话：508-647-7000\\
传真：508-647-7001\\
电子邮件：\texttt{info@mathworks.com}\\
网站：\url{www.mathworks.com}
\end{flushleft}
